\section*{Introdução}

Nesse trabalho, será feito um estudo sobre diferentes técnicas de estimação do campo de velocidade de propagação do som na coluna d'água (SSF), utilizando informação de tempo de trânsito. De forma geral, todas as técnicas que serão apresentadas baseiam-se no esquema de minimização quadrática do erro de tempo de trânsito entre o estimado (linearizado) e o mensurado, esse sendo otimizado conjuntamente da solução da equação Eikonal de forma iterativa.

Na \cref{sec:modelo-padrão-TTT} estabeleceremos um modelo-padrão para a Tomografia de Tempo de Trânsito (TTT), que é a técnica mais comum e tradicional para resolver o problema previamente estabelecido. Na \cref{sec:svd} nós desenvolveremos uma abordagem para resolver esse mesmo problema, baseada na Decomposição em Valores Singulares (SVD), juntamente de todos os desenvolvimentos matemáticos necessários comprovando sua aplicabilidade nesse problema. Essas duas técnicas serão comparadas na \cref{sec:simulation_and_results}, em ambientes simulados que emulam situações reais, com dados obtidos de \ingles{databases} e modelos conhecidos para cálculo de perfis de velocidade.

Posteriormente, nas \cref{sec:modelos_adaptados_ttt,sec:modelos_alternativos_est-ssf} serão apresentados e desenvolvidos métodos alternativos que utilizem informação de tempo de trânsito para estimação do SSF, que ou adaptem a TTT tradicional para incluir informações ou regularizações extras, ou adaptem-na para estruturar a minimização de uma maneira diferente. Os resultados para esses métodos adaptados são apresentados na \cref{sec:resultados_metodos_alternativos}.

Na \cref{sec:mais_simulacoes_raso} nós retornamos ao TTT e ao SVD apresentados respectivamente nas \cref{sec:modelo-padrão-TTT,sec:svd} (por serem os métodos estudados que não dependem de informação \ingles{a priori}, essa sendo possivelmente escassa), e avaliamos eles para uma quantidade muito maior de configurações dos dispositivos (fontes e receptores) a serem alocados na coluna d'água, com o intuito de explorar a performance das duas técnicas escolhidas em configurações que possam melhor se alinhar com situações práticas ou desejáveis, possibilitando melhor entender a relação das configurações de dispositivos no ambiente com as performances alcançáveis.

Esse estudo é posteriormente expandido na \cref{sec:mais_simulacoes_profundo}, onde além de estudar uma grande gama de configurações de dispositivos, nós também modificamos o ambiente para um com uma profundidade muito maior. O objetivo é o mesmo que da \cref{sec:mais_simulacoes_raso}, contudo agora em um ambiente diferente, cujos campos a serem estimados tem uma variação menos bem comportada em relação à profundidade.